现实数据常使 \(Ax=b\) 没有精确解：方程多于未知数，测量又带有误差，\(b\) 往往不在 \(A\) 的列空间中。问题没有因此结束，而是改变了形状：在所有可以写成 \(Ax\) 的向量里，哪一个离 \(b\) 最近？

内积给出长度和垂直的语言，正交分解把观测拆成可由模型解释的部分与残差，投影给出最近点，QR 分解则把同一几何原则变成稳定计算。四篇围绕同一个无解方程逐步推进。

以三个点 \((1,1),(2,2),(3,2)\) 拟合直线 \(y=C+Dt\) 为例：前两式给出 \(D=1,\;C=0\)，第三式却给出 \(3=2\)。矩阵形式下 \(A=\begin{pmatrix}1&1\\1&2\\1&3\end{pmatrix},\;b=(1,2,2)^{\trans}\)。列空间是 \(\R^3\) 中的平面，最小二乘在这个平面上找到 \(\hat b=(7/6,5/3,13/6)^{\trans}\)，对应 \(\hat C=2/3,\;\hat D=1/2\)，残差 \(\hat r=(-1/6,1/3,-1/6)^{\trans}\) 垂直于该平面。

\subsection{核心主线}

\begin{enumerate}
\tightlist
\item \hyperref[sec:03-Linear-Algebra/03-Orthogonality/01]{``正交与正交补''}：建立几何语言，并看见四个基本子空间为什么成对正交；
\item \hyperref[sec:03-Linear-Algebra/03-Orthogonality/02]{``投影与最近点''}：从直线上的最近点推到一般子空间，得到残差正交条件；
\item \hyperref[sec:03-Linear-Algebra/03-Orthogonality/03]{``最小二乘''}：把 (b) 投影到列空间，并区分最佳拟合存在与参数唯一；
\item \hyperref[sec:03-Linear-Algebra/03-Orthogonality/04]{``正交基与 QR 分解''}：选取让投影变简单的坐标，并避免显式形成 \(A^{\trans}A\)。
\end{enumerate}

\subsection{怎么读这一章}

前三篇是首次阅读的概念主干：先建立正交分解，再把最近点条件翻译成投影和最小二乘。若只想先理解几何，可以暂缓第四篇的 Gram--Schmidt 与 QR 算法细节；需要实际计算或判断数值可信度时，第四篇不可跳过。正规方程用于推导很方便，但在数值计算中不推荐直接形成 \(A^{\trans}A\)；秩亏时还必须区分最佳拟合值是否唯一与参数是否唯一。

\subsection{阅读前需要什么}

读者应熟悉向量、矩阵乘法、列空间、零空间、线性无关和基；需要时可定向回到\hyperref[ch:018]{``向量空间、基与线性变换''}。正文会恢复关键背景，不要求先记住投影矩阵公式。

读完本章后，应能从``残差与可行方向正交''重新推导投影和正规方程，并解释秩亏时发生什么、为什么实际计算使用 QR 或 SVD 而非直接求 \((A^{\trans}A)^{-1}\)。公式

\[
P=A(A^{\trans}A)^{-1}A^{\trans}
\]

只是这一几何原则在满列秩坐标中的表达。
